% ==========================================================================
% theory/realizer_law_theorem.tex  --  replacement for Conjecture 5.3 of MRT001 v0.5
%
% Drop-in replacement for the block of manuscript/main.tex that runs from
%   \begin{conjecture}\label{conj:realizers}
% to the end of the paragraph that starts "\noindent The conjecture is stated, not proved."
% (just before the E5e/E5f tables).  It uses only macros already defined in main.tex
% (\incomp, \status, \R, the theorem environments, \HasRlaw and the \Rlaw... /
% \EfiveE... macros of numbers.tex).
%
% Other edits needed in main.tex (NOT made here; main.tex must not be touched by this agent):
%  * every \ref{conj:realizers} -> \ref{thm:realizers} (caption of Table e5f, claims table);
%  * claims table row "realizer count = 2^N ..." : status "conjectural" -> "proved here",
%    and add a row for Proposition~\ref{prop:exceptions} (exception rate 5/n + O(n^{-2})), "proved here";
%  * the sentence of Sec. 5 "which suggests the following." -> "and this is a theorem.";
%  * in the E5e paragraph, "the product formula ... gives 2*2^N orientations" is now
%    Lemma~\ref{lem:inflation} plus Gallai's theorem;
%  * add to refs.bib the entries in theory/realizer_law_refs.bib (gallai1967, albert2005,
%    wolfowitz1944); kaplansky1945, golumbic1980 already exist.
% ==========================================================================

\paragraph{The realizer law.}
Sort the points of a configuration by $u$ and let $\pi(i)$ be the $v$-rank of the $i$-th point. For independent uniform points $u$ and $v$ are independent and continuous, so $\pi$ is uniform on the symmetric group $S_n$, and $p_i\prec p_j$ iff $i<j$ and $\pi(i)<\pi(j)$. The unrelated pairs are the inversions of $\pi$, so the incomparability graph of the order is the permutation graph $G_\pi$ of $\pi$ (vertices $1,\dots,n$, an edge $ij$ iff $(i-j)(\pi(i)-\pi(j))<0$), and by Proposition~\ref{prop:causnonunique} the number of realizers modulo the swap is $R(\pi)=1$ if $\pi$ is the identity and $R(\pi)=t(G_\pi)/2$ otherwise, where $t(G)$ is the number of transitive orientations of $G$. A pair of points adjacent in both coordinate orders with opposite directions is an index $i$ with $\pi(i+1)=\pi(i)-1$ (a \emph{descending succession}); $N(\pi)$ denotes their number. An \emph{interval} of $\pi$ is a set of consecutive positions whose values are consecutive integers; it is trivial if it has $1$ or $n$ elements, and $\pi$ is \emph{simple} if all its intervals are trivial. A \emph{module} of a graph is a set $M$ of vertices such that every vertex outside $M$ is adjacent to all or to none of $M$; it is trivial if $|M|\le1$ or $M$ is the whole vertex set, and the graph is \emph{prime} if all its modules are trivial. For $\sigma\in S_m$ and permutations $\alpha_1,\dots,\alpha_m$, the inflation $\sigma[\alpha_1,\dots,\alpha_m]$ replaces the $j$-th entry of $\sigma$ by a block of consecutive positions with consecutive values in the pattern $\alpha_j$, the blocks being ordered by value as the entries of $\sigma$. We use one classical result from the theory of comparability graphs.

\begin{theorem}[Gallai~\citep{gallai1967}]\label{thm:gallai}
A prime comparability graph has exactly two transitive orientations, one the reverse of the other.
\end{theorem}

\noindent Permutation graphs are comparability graphs ($G_\pi$ is the comparability graph of the order $i<j$, $\pi(i)>\pi(j)$), so Theorem~\ref{thm:gallai} applies to them. In a transitive orientation we write $x\to y$; the basic forcing rule~\citep{gallai1967,golumbic1980} is that if $x$ is adjacent to $a$ and $b$ and $a$, $b$ are not adjacent, then $x\to a$ iff $x\to b$, since $a\to x\to b$ would force the edge $ab$.

\begin{lemma}[Intervals and modules]\label{lem:simpleprime}
Every interval of $\pi$ is a module of $G_\pi$. Conversely, if $n\ge3$ and $G_\pi$ has a module $M$ with $2\le|M|\le n-1$, then $\pi$ has an interval with between $2$ and $n-1$ elements. Hence, for $n\ge3$, $\pi$ is simple iff $G_\pi$ is prime, and no permutation of length $3$ is simple. \status{proved here}
\end{lemma}
\begin{proof}
A point outside an interval lies on one side of it in position and on one side of it in value, so it forms an inversion with all of its points or with none. For the converse take such an $M$ of minimal size. If $G_\pi[M]$ and its complement are both connected, $M$ is an interval: let $x\notin M$ lie strictly between the smallest and the largest position of $M$; if $x$ is adjacent to no point of $M$, the points of $M$ to its left lie below it and those to its right above it, so there is no edge of $G_\pi[M]$ between the two (non-empty) groups; if $x$ is adjacent to all of $M$, every pair formed by a left and a right point of $M$ is an edge, so the complement of $G_\pi[M]$ is disconnected; either way a contradiction. The same argument applied to values shows that the values of $M$ are consecutive. If $G_\pi[M]$ is disconnected, each of its components is a module of $G_\pi$, so by minimality all are single points, $M$ is independent, every pair in $M$ is a module, and $|M|=2$; if the complement of $G_\pi[M]$ is disconnected the same argument with co-components gives $|M|=2$. Since no graph on two or three vertices is both connected and co-connected, it remains to treat a module $M=\{a,b\}$, $a<b$. If $\pi(a)<\pi(b)$ (no edge), a point $x$ with $a<x<b$ adjacent to $a$ and $b$ would need $\pi(x)<\pi(a)$ and $\pi(x)>\pi(b)$, so it is adjacent to neither and $\pi(a)<\pi(x)<\pi(b)$; and a point outside $[a,b]$ with value between $\pi(a)$ and $\pi(b)$ would be adjacent to exactly one of $a,b$. So $[a,b]$ is an interval; if it is all of $[1,n]$ then $\pi(1)=1$ and $[2,n]$ is an interval with $n-1\ge2$ elements. The case $\pi(a)>\pi(b)$ is symmetric (with $\pi(1)=n$ in the extreme case). Finally, every permutation of length $3$ has two consecutive positions with consecutive values.
\end{proof}

\begin{lemma}[Inflation of a simple permutation]\label{lem:inflation}
If $\sigma\in S_m$ is simple, $m\ge4$, and $\pi=\sigma[\alpha_1,\dots,\alpha_m]$, then $t(G_\pi)=2\prod_{j=1}^m t(G_{\alpha_j})$. \status{proved here, using Theorem~\ref{thm:gallai}}
\end{lemma}
\begin{proof}
The blocks $B_1,\dots,B_m$ are intervals, hence modules (Lemma~\ref{lem:simpleprime}), and the quotient graph is $G_\sigma$, which is prime by Lemma~\ref{lem:simpleprime}; in particular it has no universal and no isolated vertex, and its only modules with at least two elements are the whole vertex set. Fix a transitive orientation and a block $B$, and let $S$ be the set of vertices $y\notin B$ adjacent to $B$ for which there are $b_1,b_2\in B$ with $b_1\to y\to b_2$. If $y\in S$ and $z\notin B$ is adjacent to $y$ but not to $B$, the forcing rule at $y$ gives $z\to y$ (from $b_1\to y$) and $y\to z$ (from $y\to b_2$), which is absurd; so the neighbours of $y$ outside $B$ are adjacent to $B$. If $y\in S$ and $w\notin B$ is adjacent to $B$ but not to $y$, the forcing rule at $b_1$ and at $b_2$ gives $b_1\to w\to b_2$, so $w\in S$. Hence $B\cup S$ is a module. If $S\ne\emptyset$, the union of the blocks that meet $B\cup S$ is again a module (adding a block that meets a module gives a module) made of at least two blocks, so it corresponds to a module of $G_\sigma$ with at least two vertices and is everything; then a vertex $z\notin B\cup S$ would make its block a universal or isolated vertex of $G_\sigma$, since $z$ sees all of $B\cup S$ in the same way and $B\cup S$ meets every block, and $B\cup S$ equal to everything would make $B$ universal in $G_\sigma$. Both are impossible, so $S=\emptyset$: for every $y\notin B$ adjacent to $B$ all edges between $y$ and $B$ point the same way, and therefore all edges between two adjacent blocks point the same way. A transitive orientation of $G_\pi$ is thus determined by its restriction to one representative per block, a transitive orientation of $G_\sigma$, and by its restrictions to the blocks. Conversely, lifting a transitive orientation of $G_\sigma$ to all edges between blocks and choosing any transitive orientation inside each block gives a transitive orientation: for a path $a\to b\to c$, if the three points lie in different blocks or two consecutive ones share a block, transitivity follows from that of the quotient, $a$ and $c$ cannot share a block with $b$ outside it (the quotient would contain $[a]\to[b]\to[a]$), and three points in one block are handled inside the block. With $t(G_\sigma)=2$ (Theorem~\ref{thm:gallai}) the count follows.
\end{proof}

\begin{theorem}[Realizer law]\label{thm:realizers}
Let $R_n$ be the number of realizers modulo the swap of the causal order of $n$ independent uniform points of the $1{+}1$ diamond, and $N_n$ the number of pairs of points adjacent in both coordinate orders with opposite directions. For $n\ge20$:
\begin{enumerate}[label=(\alph*),leftmargin=2em]
\item $\Pr(R_n\ne2^{N_n})\le 10/n+171/n^2$;
\item $\mathbb E[N_n(N_n-1)\cdots(N_n-r+1)]=1-r/n$ for $0\le r\le n-1$, and $d_{\mathrm{TV}}(N_n,\mathrm{Po}(1))\le e^2/n+(2^n+1)/(2\,n!)$;
\item $d_{\mathrm{TV}}(R_n,2^{Z})\le(10+e^2)/n+172/n^2$, where $Z\sim\mathrm{Po}(1)$.
\end{enumerate}
In particular the fractions of samples with one, two and at most eight realizers tend to $e^{-1}$, $e^{-1}$ and $\tfrac83e^{-1}$, each with error $O(1/n)$. \status{proved here}
\end{theorem}
\begin{proof}
\emph{Step 1 (deterministic).} Let $n\ge5$ and suppose that $\pi$ has no interval with $k$ elements for $3\le k\le n-1$. Two overlapping successions (descending or ascending) would form an interval with three elements, so the successions are disjoint. Contracting each of them to a point gives $\sigma\in S_m$ with $m\ge\lceil n/2\rceil\ge3$ and $\pi=\sigma[\alpha_1,\dots,\alpha_m]$, $\alpha_j\in\{1,12,21\}$. An interval of $\sigma$ with between $2$ and $m-1$ elements would expand to an interval of $\pi$ with between $2$ and $n-1$ elements, which is excluded if it has at least three elements and would be an uncontracted succession if it has two; so $\sigma$ is simple, and $m\ge4$ by Lemma~\ref{lem:simpleprime}. Lemma~\ref{lem:inflation} with $t(G_{1})=t(G_{12})=1$ and $t(G_{21})=t(K_2)=2$ gives $t(G_\pi)=2\cdot2^{N(\pi)}$, and since $\pi$ is not the identity, $R(\pi)=2^{N(\pi)}$.

\emph{Step 2 (first moment).} Let $I_k$ be the number of intervals of $\pi$ with $k$ elements. Choosing the first position, the smallest value and the arrangements inside and outside gives $\mathbb E I_k=(n-k+1)^2k!\,(n-k)!/n!$. Thus $\mathbb E I_3=6(n-2)/(n(n-1))\le6/n$, $\mathbb E I_{n-1}=4/n$, $\mathbb E I_{n-2}=18/(n(n-1))$, $\mathbb E I_{n-3}=96/(n(n-1)(n-2))$, $\mathbb E I_4=24(n-3)/(n(n-1)(n-2))$, $\mathbb E I_{n-4}=600/(n(n-1)(n-2)(n-3))$, and for $5\le k\le n-5$, $\mathbb E I_k\le(n-4)^2/\binom n5$, so these $n-9$ terms add up to at most $120(n-4)(n-9)/(n(n-1)(n-2)(n-3))\le120/n^2$. For $n\ge20$ the five middle terms are at most $(24+19+6+2+120)/n^2$, and by Step~1 and the union bound
\[
\Pr(R_n\ne2^{N_n})\le\sum_{k=3}^{n-1}\mathbb E I_k\le\frac{10}{n}+\frac{171}{n^2},
\]
which is (a); the sum equals $10/n+O(n^{-2})$.

\emph{Step 3 (factorial moments).} Let $A$ be a set of $r$ indices in $\{1,\dots,n-1\}$. The event that $\pi(i+1)=\pi(i)-1$ for all $i\in A$ says that each maximal run of consecutive indices of $A$, $i,\dots,i+l-1$, occupies positions $i,\dots,i+l$ with consecutive decreasing values; contracting each run to a point is a bijection onto $S_{n-r}$, so the event has probability $(n-r)!/n!$ whatever $A$ is. Hence $\mathbb E[N_n(N_n-1)\cdots(N_n-r+1)]=r!\binom{n-1}r(n-r)!/n!=1-r/n$, the first part of (b). Since $N_n\le n-1$, inclusion--exclusion is exact:
\[
\Pr(N_n=j)=\frac1{j!}\sum_{s=0}^{n-1-j}\frac{(-1)^s}{s!}\Bigl(1-\frac{j+s}n\Bigr),
\]
so $|\Pr(N_n=j)-e^{-1}/j!|\le1/(j!\,(n-j)!)+e(j+1)/(n\,j!)$; summing over $j$ and adding the Poisson mass beyond $n-1$ gives $\sum_j|\Pr(N_n=j)-e^{-1}/j!|\le2e^2/n+(2^n+1)/n!$, the second part of (b). This is the classical Poisson limit of successions in a random permutation~\citep{wolfowitz1944,kaplansky1945}, stated there for ascending successions; reversing the order of the positions exchanges the two kinds and preserves the uniform law.

\emph{Step 4.} $d_{\mathrm{TV}}(R_n,2^Z)\le\Pr(R_n\ne2^{N_n})+d_{\mathrm{TV}}(2^{N_n},2^Z)$, and the last term equals $d_{\mathrm{TV}}(N_n,Z)$ because $j\mapsto2^j$ is injective; for $n\ge20$, $(2^n+1)/(2\,n!)\le1/n^2$. The limits follow from $\Pr(Z=0)=\Pr(Z=1)=e^{-1}$ and $\Pr(Z\le3)=\tfrac83e^{-1}$.
\end{proof}

\begin{proposition}[Where the exceptions come from]\label{prop:exceptions}
$\Pr(R_n\ne2^{N_n})=3(n-2)/(n(n-1))+2/n+O(n^{-2})=5/n+O(n^{-2})$. Up to an event of probability $O(n^{-2})$, an exception is a configuration with a single interval of three elements in one of the patterns $321$ (then $R_n=\tfrac32\,2^{N_n}$), $231$ or $312$ (then $R_n=2\cdot2^{N_n}$), or with $\pi(1)=n$ or $\pi(n)=1$, i.e.\ a point related to no other one and extreme in both coordinate orders (then $R_n=2\cdot2^{N_n}$). \status{proved here}
\end{proposition}
\begin{proof}
By Step~2 the probability that $\pi$ has an interval with between $4$ and $n-2$ elements is $O(n^{-2})$, and so is the probability of two intervals with $3$ or $n-1$ elements: two disjoint three-element intervals have expected number at most $18n^2(n-4)!/n!$, two overlapping ones form an interval with four or five elements, a three-element interval together with $\pi(1)\in\{1,n\}$ or $\pi(n)\in\{1,n\}$ has probability at most $24/n^2+12/(n(n-1))$, and $\pi(1),\pi(n)\in\{1,n\}$ has probability $2/(n(n-1))$. Outside this event, if $\pi$ has no interval with between $3$ and $n-1$ elements, $R_n=2^{N_n}$ by Step~1. If it has exactly one, $J$, with three elements, contracting $J$ and the successions disjoint from it gives a simple permutation as in Step~1, and Lemma~\ref{lem:inflation} gives $R_n/2^{N_n}=t(G_p)/2^{N(p)}$ for the pattern $p$ of $J$: this is $1$ for $123$, $132$, $213$ ($t=1,2,2$; $N(p)=0,1,1$), $6/4=\tfrac32$ for $321$ ($G_p=K_3$, $t=3!$) and $2$ for $231$ and $312$ ($G_p$ a path on three vertices, $t=2$, $N(p)=0$). If $J$ has $n-1$ elements, $\pi$ is $1\oplus\pi'$, $\pi'\oplus1$, $1\ominus\pi'$ or $\pi'\ominus1$ (by $\oplus$ and $\ominus$ we mean the direct and skew sums, i.e.\ $\pi(1)=1$, $\pi(n)=n$, $\pi(1)=n$, $\pi(n)=1$), where $\pi'$ satisfies the hypothesis of Step~1 and is neither a direct nor a skew sum (otherwise $\pi$ would have a second interval with at least three elements); then $G_\pi$ is $G_{\pi'}$ plus an isolated vertex in the first two cases and $G_{\pi'}$ plus a vertex adjacent to all others in the last two, and $N(\pi)=N(\pi')$. In the last two cases the forcing rule applied to the non-edges of the co-connected graph $G_{\pi'}$ makes the extra vertex a source or a sink, so $t(G_\pi)=2\,t(G_{\pi'})$ and $R_n=2\cdot2^{N_n}$; in the first two $R_n=2^{N_n}$. Each three-element pattern has expected number of occurrences $(n-2)/(n(n-1))$ and $\Pr(\pi(1)=n)=\Pr(\pi(n)=1)=1/n$, so the upper bound follows, and Bonferroni's inequality with the $O(n^{-2})$ bounds on intersections gives the matching lower bound.
\end{proof}

\noindent Lemmas~\ref{lem:simpleprime} and~\ref{lem:inflation}, together with the decomposition of every permutation with at least two elements as a direct sum, a skew sum, or an inflation of a simple permutation with at least four elements~\citep{albert2005}, give by induction the exact count for every $\pi\ne\mathrm{id}$: $R(\pi)$ is half the product, over the nodes of the substitution decomposition tree of $\pi$, of $2$ for a simple node, $k!$ for a skew-sum node with $k$ children and $1$ for a direct-sum node. This is Gallai's product formula~\citep{gallai1967} for the incomparability graph, in which skew sums are series nodes and direct sums parallel nodes. Thus a descending succession is a series module $K_2$ (factor $2$) and an ascending one is a parallel module (factor $1$); factors $3$ and higher come from a skew-sum root with three children, e.g.\ $\pi(1)=n$ together with $\pi(n)=1$, or from combinations of the cases of Proposition~\ref{prop:exceptions}, all with probability $O(n^{-2})$. \status{classical (Gallai; Albert--Atkinson), checked against brute force for all $\pi$ with $n\le7$}
\if\HasRlaw1 The predictions \RlawPoissonOne, \RlawPoissonOne{} and \RlawPoissonLeThree{} of Theorem~\ref{thm:realizers} are consistent with Table~\ref{tab:e5e}, where the three fractions at $n=\EfiveEnmax$ are \EfiveEuniqueNmax, \EfiveEtwoNmax{} and \EfiveEleEightNmax, and with the independent samples of E5f (Table~\ref{tab:e5f}): the exact count equalled $2^N$ in a fraction \RlawEqualFracNmin{} of the samples at $n=\RlawNmin$, in \RlawEqualBig{} of the \RlawSamplesBig{} samples with $n\ge\RlawBigNmin$, and in \RlawEqualNmax{} of \RlawSamplesNmax{} at $n=\RlawNmax$. The exceptions are counts that differ from $2^N$ by a factor $3/2$, $2$ or $3$, as described by Proposition~\ref{prop:exceptions}, whose rate $5/n$ ($0.25$ at $n=20$, $0.10$ at $n=50$) matches the observed fractions. \fi Thus the realizer of the causal order of a uniform sample is unique with probability tending to $e^{-1}$, and in general it is determined up to $N_n$ independent binary choices, one per pair of points adjacent in both coordinate orders with opposite directions. This is the finite analogue of the continuum theorems, and it differs from the Euclidean ordinal case, where the kernel class always has non-empty interior.
